% =====================================================================
%  c00 —— 标题页 / 摘要 / 译者说明（Lead 独占）
%  原文 p.1：标题、作者、摘要、目录
% =====================================================================
% [原文 p.1]
\pagenumbering{roman}

\begin{titlepage}
\centering
\vspace*{2.5cm}
{\LARGE\bfseries 规范场论中的进阶专题：\par}
\vspace{0.3cm}
{\LARGE\bfseries Higgs 丛的数学与物理\par}
\vspace{1.6cm}
{\large Laura P. Schaposnik\par}
\vspace{0.4cm}
{\normalsize University of Illinois at Chicago\par}
\vspace{2.2cm}
{\large 中文译注本\par}
\vspace{0.3cm}
{\normalsize（据 arXiv:1907.09800v1，2019 年 7 月 23 日；PCMI 2019 研究生暑期学校讲义）\par}
\vfill
{\small 本译本由 AI 协作完成：正文逐段对照原文页图翻译，公式逐符号核对；\par}
{\small 全部译者批注以 \texttt{\% [Tutor Note: ...]} 形式存放于源码，未进入排版。\par}
\vspace{0.8cm}
{\small \today\par}
\end{titlepage}

% ---------- 摘要 ----------
\thispagestyle{plain}
\begin{center}
{\large\bfseries 摘要}
\end{center}

\noindent
本讲义是作者在 2019 年帕克城数学研究所（Park City Mathematics Institute, PCMI）与高等研究院（Institute for Advanced Study）联合举办的研究生暑期学校所授短课程（mini-course）的阅读材料。我们首先介绍 Higgs 丛及其主要性质（第一讲），随后讨论 Hitchin 纤维化及其不同用途（第二讲）。课程的后半部分致力于研究复 Higgs 丛模空间中的各类子空间（即「膜」，brane），它们在平坦联络与表示论中的表现（第三讲），以及这些膜之间的对应关系（第四讲）。

\vspace{0.4cm}
\noindent
{\small\textbf{2010 数学主题分类}：Primary ????；Secondary ????（原文如此，arXiv 版未填写——见译者说明）\par}
\noindent
{\small\textbf{关键词与短语}：Park City Mathematics Institute。\par}

% ---------- 译者说明 ----------
\vspace{0.8cm}
\begin{center}
{\large\bfseries 译者说明}
\end{center}

\noindent
本译注本以 arXiv:1907.09800v1（53 页）为底本，按「原文页码」逐页对照翻译，原文印刷页码与本译本页面的对应关系见交付包 \texttt{README.md}。体例要点如下：

\begin{itemize}[leftmargin=2em,itemsep=2pt]
\item 所有编号公式、定理、定义、例、注、猜想、问题、图、表的编号\textbf{一律照抄原文}；原文在每一节内使用统一计数器（例如 $\S2$ 中「例 2.10」之后紧接的就是式 (2.12)），本译本保持这一编号体系，绝不自行重排。
\item 原文的\textbf{数学记号逐一按页图核对}——文字层中的数学抽取不可信，因此每个公式都以 200\,dpi 页图为唯一真值。
\item 原文\textbf{笔误照原样排}，仅在批注中指出并给出正确形式；疑似笔误一律先裁 500\,dpi 以上原图复核。
\item 原文参考文献（p.49--53）按学术惯例\textbf{不译}，以英文原貌转录，编号与原文一致。
\item 译者批注（Tutor Note）不进入排版，单独汇总于 \texttt{译注汇总.md}。
\end{itemize}

\clearpage
\tableofcontents
\clearpage
